\documentclass[../main]{subfiles}
\begin{document}
\chapter*{Introducción}
Esta propuesta formativa plantea la apropiación de las características
constructivas y los fenómenos electromagnéticos que dan vida a las
\emph{Máquinas
Eléctricas} y al funcionamiento de la maquinaria eléctrica más
empleada en el sector
técnico, ya sea como motores o generadores eléctricos. Se analizan los motores
asincrónicos trifásicos, con sus principios de funcionamiento, conceptualizando
comportamientos energéticos de rotores y representándolos a través de diagramas
vectoriales y circuitos equivalentes. Como saberes eléctricos se
destacan las curvas
que hacen referencia a los arranques de los motores eléctricos, ya sea de forma
directa o estrella-triángulo, dependiendo de la potencia eléctrica
instalada. También se
analizarán los tipos de arranque resistivo, reactivo, con empleo
diverso de capacitores,
etc., al observar el funcionamiento de las máquinas monofásicas en particular.

Las teorías de fundamento desarrolladas con anterioridad en Circuitos
Eléctricos II, Máquinas Eléctricas I y Electrotecnia II propician la
comprensión de los
campos electromagnéticos alternativos y rodantes, con sus fuerzas
electromotrices
respectivas. Se desarrollan proyectos articulados con los espacios curriculares
de Seguridad e Higiene en el Trabajo, Laboratorio de Mediciones
Eléctricas II, Obras
Eléctricas, Electrónica Industrial y Operaciones de Taller de
Electricidad II. Es
importante la permanente articulación con el Laboratorio de Ensayos
Eléctricos, para
efectuar las prácticas con el instrumental correspondiente, tanto de los motores
asincrónicos trifásicos, como de las máquinas monofásicas, motores de
denominación
universal y todas las máquinas que son objeto de estudio de este
espacio curricular.

La presente propuesta formativa integra teoría y práctica, conforme lo
establece el artículo 35 de la Resolución CFE 229/14.

\begin{enumerate}
  \item FUNDAMENTOS DE MÁQUINAS MONOFÁSICAS
    \begin{itemize}
      \item Reconocer los principios básicos que fundamentan el
        funcionamiento de una máquina eléctrica con energía alterna.
        \begin{itemize}
          \item Explicitación de las leyes electromagnéticas presentes en
            una máquina eléctrica.
          \item Interpretación de los campos electromagnéticos
            originados por corriente alterna.
          \item Distinción de la frecuencia y la velocidad de sincronismo
            en una máquina eléctrica.
          \item Caracterización de la electrodinámica en máquinas
            industriales.
        \end{itemize}

      \item Interpretar las teorías electrodinámicas eléctricas en
        una máquina rotativa.
    \end{itemize}
  \item MOTORES ASINCRÓNICOS TRIFÁSICOS EN DIVERSOS ENSAYOS ELÉCTRICOS
    \begin{itemize}
      \item Conocer la naturaleza
        eléctrica de las máquinas
        asincrónicas.
      \item Analizar críticamente los
        ensayos eléctricos
        característicos de un motor
        asincrónico que
        fundamentan los distintos
        fenómenos energéticos.

    \end{itemize}
  \item ENSAYOS EN MÁQUINAS SINCRÓNICAS
    \begin{itemize}
      \item Comprender el
        funcionamiento de las
        partes constitutivas de una
        máquina sincrónica.
      \item Valorar la utilidad industrial
        de una máquina sincrónica
        considerando su
        importancia tecnológica
        eléctrica.
    \end{itemize}
\end{enumerate}

\subsection*{Aprendizajes Prioritarios}
\begin{itemize}
  \item Explicitación de las leyes electromagnéticas presentes en una
    máquina eléctrica.
  \item Interpretación de los campos electromagnéticos originados por
    corriente alterna.
  \item Distinción de la frecuencia y la velocidad de sincronismo en
    una máquina eléctrica.
  \item Caracterización de la electrodinámica en máquinas industriales.
  \item Interpretación de los ángulos geométricos y eléctricos.
  \item Comprensión de la fuerza electromotriz inducida en máquinas rotativas.
  \item Reconocimiento de los factores de arrollamiento y de paso.
  \item Análisis enmarcado en la tecnológica eléctrica de los usos y
    características de los motores paso a paso.
  \item Interpretación de los aspectos constructivos de un motor asincrónico.
  \item Reconocimiento del principio de funcionamiento de una máquina asíncrona
    y su resbalamiento.
  \item Identificación del motor eléctrico como transformador
    energético en su empleo técnico.
  \item Análisis del comportamiento eléctrico con rotor detenido y en marcha.
  \item Elaboración de los diagramas vectoriales.
  \item Interpretación gráfica de los circuitos equivalentes
    técnicamente eléctricos.
  \item Análisis de las curvas de arranque, cuplas y balance de potencias.
  \item Distinción de formas de arranque de un motor trifásico.
  \item Reconocimiento de la máquina asíncrona como autotransformador.
  \item Contrastación de los distintos tipos de frenado eléctrico en sus
    distintas concepciones tecnológicas.
  \item Identificación de los principios de funcionamiento del generador
    y del motor eléctrico.
  \item Descripción de las máquinas sincrónicas en vacío y en carga.
  \item Interpretación de las partes constitutivas de las máquinas sincrónicas.
  \item Caracterización de la reacción de inducido.
  \item Diferenciación en la regulación de una máquina sincrónica.
  \item Selección de la metodología de sincronización de generadores
    sincrónicos.
  \item Reconocimiento tecnológico de una máquina sincrónica en su rol
    de generador o motor eléctrico
\end{itemize}
\tableofcontents
\end{document}
